% ----------------------
\subsection{Section 88 (27--28 December 1832; 3 January 1833)}
% ----------------------
	\label{DC88}
	
	% -----
	\subsubsection{Historical Background}
	% -----
		\label{DC88-HB}
		
		% --
		\paragraph{JSP}
		% --
		
			%
			\subparagraph{Verses 1--126}
			%
			
				``JS dictated a lengthy revelation at a conference of high priests in Kirtland, Ohio, on 27--28 December 1832. The revelation's heading, which was probably provided by Frederick G. Williams, states that the revelation was addressed to the `first Elders' of the church. The text of the revelation describes its audience as those who had congregated at the conference in Kirtland `to receive his [God's] will concerning you.' Although in later years the term `first Elders' generally referred to the leading elders of the church, here it appears to have a less hierarchical meaning, equating the elders to whom the revelation was addressed with `the first \sout{Elders} labourers, in this last kingdom' who were referenced in a parable presented in this revelation. A later JS history emphasizes that the revelation came two days after a revelation describing an outbreak of wars and slave rebellions that would precede Christ's second coming.

				``JS called this revelation `the Olieve leaf which we have plucked from the tree of Paradise' and `the Lords message of peace to us.' Perhaps JS described the revelation in this way because it offset the stark apocalyptic imagery of the 25 December revelation or perhaps because he saw its messages regarding the conduct of church members and the need for unity as a way to heal ongoing difficulties with Missouri church leaders. Like the 25 December prophecy of war, the 27--28 December revelation discusses eschatological events, but interspersed throughout the revelation are explanations of the requirements to enter the celestial, terrestrial, and telestial kingdoms in the life to come and an exposition on light and its relation to Jesus Christ.

				``Much like the first chapter of the book of John (which JS revised in late 1831 or early 1832 as part of his Bible revision), the first part of this revelation connects Christ with light and the creative process. This explanation expanded on ideas expressed in earlier revelations. Revelations in 1829, for example, generally used the concept of light to represent Jesus Christ. By 1831, revelations were also using light as a metaphor for the gospel and as a more abstract representation of truth and knowledge. The 27--28 December revelation brings such ideas together by explaining that Christ's light, which the revelation defines as truth and knowledge, is in all things, is the power by which they were created, and is the law governing them. Such concepts were not entirely novel; in the 1700s, Swedish theologian Emanuel Swedenborg, for example, argued that `the light which proceeds from the Lord as a sun is Divine Truth, from which the angels derive all their wisdom and intelligence,' but this revelation goes further in its connection of light to the creative and governing processes.

				``While it explored theological themes, the revelation also issued concrete directives, instructing the elders to sanctify themselves at a `solemn assembly,' to construct a house of God, and to be taught there in both spiritual and temporal matters before embarking on their missions to the Gentiles `for the last time.' These instructions came in response to specific prayers that God show `his will . . . concerning the upbuilding of Zion,' which suggests that the revelation would apply only to church members in Missouri. Saints in Kirtland, Ohio, however, took the direction as a call to action. Just two weeks after JS dictated this revelation, he informed church leaders in Missouri that the revelation provided a commandment from God `to build an house of God, \& establish a school for the Prophets' in Kirtland. Regarding the direction as `the word of the Lord to us,' JS and the Saints in Kirtland promptly began to organize the `school for the Prophets.' A revelation dictated less than a week after the 27--28 December revelation, which would later become associated with it, provided more instructions on establishing the school. Over the next several months, church leaders in Kirtland took steps to construct a schoolhouse for `the Elders who should come in to receive ther education for the ministry' and broke ground for the building they called the House of the Lord.

				``Many Saints focused more on the revelation's immediate directives than on its metaphysical aspects. Samuel Smith, for example, wrote in his journal that the revelation instructed `the first labourers in this last vinyard' to `call a sollem assembly' where they could `sanctify themselves \& wash their hands \& feet for a testimony' against an unbelieving generation. He also highlighted the revelation's requirement to `appoint a teacher among' the elders so that they could obtain `knowledge of countries \& languages.' Nowhere in his journal did Samuel Smith refer to the eschatology of the revelation or its other doctrinal points. Likewise, when William W. Phelps printed part of the revelation in the February 1833 issue of \textit{The Evening and the Morning Star}, he chose portions explaining the solemn assembly and the construction of the House of the Lord.

				``As the note at the end of the inscription indicates, Frederick G. Williams wrote this revelation as JS dictated it. The original manuscript is not extant; Williams copied the revelation into Revelation Book 2, probably between late January and late February 1833. Soon after dictating the revelation, JS transmitted it to the Saints in Missouri by enclosing a copy of the text in a letter to Phelps, explaining that its contents showed `that the Lord approves of us \& has accepted us, \& established his name in kirtland for the salvation of the nations.' This revelation was first published in its entirety on a broadside in late 1833 or early 1834. It was later combined and printed with the revelation of 3 January 1833.''

			%
			\subparagraph{Verses 127--137}
			%
			
				``According to the index in Revelation Book 2, the subject of this 3 January 1833 revelation is `instructions how to regulate the Elders school.' On 27--28 December 1832, JS dictated a revelation directing the elders of the church to establish a school so that they could be taught in both `the doctrines, of the kingdom' and in temporal matters before proclaiming the gospel to the world `for the last time.' It also instructed the Saints in Kirtland, Ohio, to build a `house of God' that, among other things, would serve as a `house of Learning.' This 3 January revelation provided information on how the school was to be conducted in the House of the Lord, including how members of the school were to act when entering it. Little is known about the circumstances of the 3 January revelation, other than that it was dictated in Kirtland, probably in Newel K. Whitney's white store, where JS and his family were staying.
				
				``Frederick G. Williams wrote the revelation as JS dictated it, but the original manuscript is not extant. Williams later copied the revelation into Revelation Book 2, probably between late January and late February 1833. On 11 January 1833, JS sent a copy of the 27--28 December 1832 revelation to William W. Phelps and other leaders in Jackson County, Missouri. It is unclear whether JS also included this revelation in that transmission, but Phelps published the revelation in the March 1833 issue of \textit{The Evening and the Morning Star}, titling it `Revelation Given Kirtland, Ohio, January 3, 1833.' Though it was presented in early manuscript sources and in \textit{The Evening and the Morning Star} as separate from the 27--28 December 1832 revelation, the 3 January 1833 revelation was combined with that earlier revelation into one text in the 1835 Doctrine and Covenants, probably because the 3 January revelation included instructions for implementing the school required by the 27--28 December revelation. The two became so closely associated that when the 3 January revelation was reprinted in \textit{Evening and Morning Star} (an 1835--1836 reprint of \textit{The Evening and the Morning Star}), it was dated 27 December 1832. Likewise, the heading for the combined revelations in the 1835 edition of the Doctrine and Covenants listed the date as 27 December 1832.''
				
		% --
		\paragraph{Kirtland Council Minute Book}
		% --
		
			[See the account of this revelation from this source \hyperref[EMS-1832-JS-27-28-DEC-1832-MIN]{here}.]
			
	% -----
	\subsubsection{Versions}
	% -----
		\label{DC88-V}
		
		See the \href{https://drive.google.com/file/d/1goAvNz8jS4R8slYfMG_db55Ty2z_vmgK/view?usp=sharing}{sections comparison} document.
	
		\begin{compactdesc}
			\item[JSP] \phantom{.}
				\begin{compactdesc}
					\item[Verses 1--126] \hyperref[EMS-1832-JS-27-28-DEC-1832]{27--28 December 1832}
					\item[Verses 127--137] \hyperref[EMS-1833-JS-3-JAN-1833]{3 January 1833}
				\end{compactdesc}
			\item[EMS] 1:9 (February 1833), 69 (verses 117--126 only)
			\item[EMS] 1:10 (March 1833), 78 (verses 127--137 only)
			\item[BC] ---
			\item[DC 1835] Section 7, 100--108
			\item[MS] 1:10 (February 1841), 245--252
			\item[TS] 5:20--21 (1 and 15 November 1844), 688-691, 704--705
			\item[DC 1844] Section 7, 130--145
			\item[MS] 14:20 (10 July 1852), 305--309
		\end{compactdesc}

	% -----
	\subsubsection{Section 88:1} % *DC88-1 
	% -----
		\label{DC88-1}

		VERILY, thus saith the Lord unto you who have assembled yourselves together to receive his will concerning you:

		\begin{framed}
		\scriptsize
			\begin{compactdesc}
				\item[JSP] Verily thus saith the Lord unto you, who have assembled yourselves together, to receive his will concerning you,
				% % \item[BC] 
				% % \item[]
			\end{compactdesc}
		\end{framed}

	% -----
	\subsubsection{Section 88:2} % *DC88-2 
	% -----
		\label{DC88-2}

		Behold, this is pleasing unto your Lord, and the angels rejoice over you; the alms of your prayers have come up into the ears of the Lord of \hyperref[SABAOTH]{Sabaoth}, and are recorded in the book%
			\footnote{%
				% \label{}%
				See \hyperref[BOOK]{Book} for this ``book'' and others.
				}
		of the names of the sanctified, even them of the celestial world.

		\begin{framed}
		\scriptsize
			\begin{compactdesc}
				\item[JSP] behold this is pleasing unto unto your lord, and the Angels rejoice over you, the \sout{prayers} alms of your prayers have come up into the ears of the Lord of sabaoth, and are recorded in the book of the names of the sanctified, even they of the celestial world,
				% % \item[BC] 
				% % \item[]
			\end{compactdesc}
		\end{framed}

	% -----
	\subsubsection{Section 88:3} % *DC88-3 
	% -----
		\label{DC88-3}

		Wherefore, I now send upon you another \hyperref[COMFORTER]{Comforter}, even upon you my friends, that it may abide in your hearts, even the \hyperref[HSOP]{Holy Spirit of promise}; which other \hyperref[COMFORTER]{Comforter} is the same that I promised unto my disciples, as is recorded in the testimony of John.%
		\footnote{%
			% \label{}%
			See \hyperref[JOHN-14-16]{John 14:16--18} and \hyperref[JOHN-14-26]{26}; \hyperref[JOHN-15-26]{John 15:26}; and \hyperref[JOHN-16-7]{John 16:7}.
			} 

		\begin{framed}
		\scriptsize
			\begin{compactdesc}
				\item[JSP] wherefore, I now send upon you another comfortor, even upon you my friends; that it may abide in your hearts, even the holy spirit of promise which other comforter, is the same, that I promised unto my deciples, as is recorded in the testamony of John,
				% % \item[BC] 
				% % \item[]
			\end{compactdesc}
		\end{framed}

	% -----
	\subsubsection{Section 88:4} % *DC88-4 
	% -----
		\label{DC88-4}

		This \hyperref[COMFORTER]{Comforter} is the promise which I give unto you of eternal life, even the glory of the celestial kingdom;

		\begin{framed}
		\scriptsize
			\begin{compactdesc}
				\item[JSP] This comfortor is the promise which I give unto you of eternal life; even the glory of the celestial kingdom,
				% % \item[BC] 
				% % \item[]
			\end{compactdesc}
		\end{framed}

	% -----
	\subsubsection{Section 88:5} % *DC88-5 
	% -----
		\label{DC88-5}

		Which glory is that of the church of the \hyperref[FIRSTBORN]{Firstborn},%
			\footnote{%
				% \label{}%
				For ``church of the Firstborn,'' see \hyperref[CHURCH-PHRASES]{here}.
				}
		even of God,%
			\footnote{%
				% \label{}%
				We're still talking about the ``glory,'' here, so the right way to read this passage is---``which glory is that of the church of the Firstborn, even [the glory] of God, the holiest of all, [and this glory comes] through Jesus Christ his Son.''
				}
		the holiest of all, through Jesus Christ his Son---%
			\footnote{%
				% \label{}%
				Note that we receive the glory, or it's transmitted to us, ``through Jesus Christ.'' We're talking here about power in the priesthood. \hyperref[DC107-3]{DC 107:3} teaches a similar principle---the priesthood is ``the Holy Priesthood, after the Order of the Son of God.'' Again, we receive it through the Son of God. Compare also \hyperref[DC84-18]{DC 84:18}.
				}

		\begin{framed}
		\scriptsize
			\begin{compactdesc}
				\item[JSP] which glory is that of the church of the first born; even of God. the holiest of all; through Jesus Christ, his son.
				% % \item[BC] 
				% % \item[]
			\end{compactdesc}
		\end{framed}

	% -----
	\subsubsection{Section 88:6} % *DC88-6 
	% -----
		\label{DC88-6}

			\footnote{%
				% \label{}%
				This verse seems like it begins an apposition to the name ``Jesus Christ his Son'' from the previous verse---the text is now describing Christ, giving some of his attributes. If this interpretation is correct, you could start this verse this way: ``This Jesus Christ is the one that ascended up on high...''
				}%
		He%
			\footnote{%
				% \label{}%
				Compare \hyperref[COLOSSIANS-1-9]{Colossians 1:9--29}, which seems like a source for some of the language in this passage.
				}
		that
			\footnote{%
				% \label{}%
				For ``ascended'' and ``descended,'' see \hyperref[EPHESIANS-4-6]{Ephesians 4:6--10}.
				}
		ascended up on high,%
			\footnote{%
				% \label{}%
				Compare \hyperref[DC132-19]{DC 132:19}.
				}
		as also he descended below%
			\footnote{%
				% \label{}%
				Compare \hyperref[DC122-8]{DC 122:8}.
				}
		all things, in that%
			\footnote{%
				% \label{}%
				The ``in that'' marks that what we are about to get is an explanation of Christ's descent below and possibly also his ascent above. He did those things \textit{in that} he comprehended all things.
				}
		he \hyperref[COMPREHEND]{comprehended}%
			\footnote{%
				% \label{}%
				I read this ``comprehend'' as drawing from \hyperref[JOHN-1-5]{John 1:5}; see that verse, and also \hyperref[KATALAMBANO]{\grl{καταλαμβάνω}}.
				}
		all things, that he might be in all and through all things,%
			\footnote{%
				% \label{}%
				See \hyperref[DC63-59]{DC 63:59}.
				}
		the light%
			\footnote{%
				% \label{}%
				See \hyperref[LUKE-2-30]{Luke 2:30--32}.
				}
		of truth;%
			\footnote{%
				% \label{}%
				For ``light of truth,'' see \hyperref[TRUTH-PHRASES]{here}.
				}

		\begin{framed}
		\scriptsize
			\begin{compactdesc}
				\item[JSP] he that assended up on high, as also he, decended below all things; in that he comprehended all things, that he might be in all, and through all things; the light of truth,
				% % \item[BC] 
				% % \item[]
			\end{compactdesc}
		\end{framed}

	% -----
	\subsubsection{Section 88:7} % *DC88-7 
	% -----
		\label{DC88-7}

		Which truth shineth. This is the light of Christ.%
			\footnote{%
				% \label{}%
				For ``light of Christ,'' see \hyperref[LIGHT-PHRASES]{here}.
				}
		As also he is in the sun, and the light of the sun,%
			\footnote{%
				% \label{}%
				I'm not sure which is the right way to read this passage---it could be, ``As also he is in the sun, and [is in] the light of the sun, and [is in] the power thereof...'', but it also could be ``As also he is in the sun, and [is] the light of the sun, and [is] the power thereof...'' I think I prefer the latter because it's easier to relate to the priesthood, but both could work and relate to the priesthood without a lot of trouble. Note that the next verse suggests that the latter interpretation is correct.
				}
		and the power thereof by which it was made.%
			\footnote{%
				% \label{}%
				See \hyperref[DC45-1]{DC 45:1}.
				}

		\begin{framed}
		\scriptsize
			\begin{compactdesc}
				\item[JSP] \sout{therefore} which <truth> shineth--- this is the light of Christ as also he is in the s$\diamond$n, and the light of the son, and the power thereof by which it was made,
				% % \item[BC] 
				% % \item[]
			\end{compactdesc}
		\end{framed}

	% -----
	\subsubsection{Section 88:8} % *DC88-8 
	% -----
		\label{DC88-8}

		As also he is in the moon, and is the light of the moon, and the power thereof by which it was made;

		\begin{framed}
		\scriptsize
			\begin{compactdesc}
				\item[JSP] as also he is in the moon, \& is, the light of the moon, and the power thereof, by which it was made,
				% % \item[BC] 
				% % \item[]
			\end{compactdesc}
		\end{framed}

	% -----
	\subsubsection{Section 88:9} % *DC88-9 
	% -----
		\label{DC88-9}

		As also the light of the stars, and the power thereof by which they were made;

		\begin{framed}
		\scriptsize
			\begin{compactdesc}
				\item[JSP] as also the light of the stars, and the power thereof; by which they were made;
				% % \item[BC] 
				% % \item[]
			\end{compactdesc}
		\end{framed}

	% -----
	\subsubsection{Section 88:10} % *DC88-10 
	% -----
		\label{DC88-10}

		And the earth also, and the power thereof, even the earth upon which you stand.

		\begin{framed}
		\scriptsize
			\begin{compactdesc}
				\item[JSP] and the earth also, \sout{and} and the power thereof, even the earth upon which you stand,
				% % \item[BC] 
				% % \item[]
			\end{compactdesc}
		\end{framed}

	% -----
	\subsubsection{Section 88:11} % *DC88-11 
	% -----
		\label{DC88-11}

			\footnote{%
				% \label{}%
				For the next three verses, compare \hyperref[DC84-45]{DC 84:45--46}.
				}%
		And the light which shineth, which giveth you light, is through him who enlighteneth your eyes,%
			\footnote{%
				\label{DC88-11-enlighteneth}%
				See \hyperref[EPHESIANS-1-17]{Ephesians 1:17--18}, ``the eyes of your understanding being enlightened.'' 
				}
		which is the same light that \hyperref[QUICKEN]{quickeneth}%
			\footnote{%
				% \label{}%
				See \hyperref[DC33-16]{DC 33:16}.
				}
		your understandings;

		\begin{framed}
		\scriptsize
			\begin{compactdesc}
				\item[JSP] and the light which now shineth; which giveth you light, is through him which enlightneth your eyes; which is the same light that quickneth your understandings,
				% % \item[BC] 
				% % \item[]
			\end{compactdesc}
		\end{framed}

	% -----
	\subsubsection{Section 88:12} % *DC88-12 
	% -----
		\label{DC88-12}

		Which light proceedeth%
			\footnote{%
				% \label{}%
				Compare \hyperref[JOHN-15-26]{John 15:26}, ``even the Spirit of truth, which proceedeth from the Father.''
				}
		forth from the presence of God%
			\footnote{%
				% \label{}%
				See \hyperref[REVELATIONNT-22-1]{Revelation 22:1}; that verse uses the word ``throne,'' which also appears here in verse 13.
				}
		to fill the immensity of space---

		\begin{framed}
		\scriptsize
			\begin{compactdesc}
				\item[JSP] which light, procedeth forth from the presence of God; to fill the emencity of space;
				% % \item[BC] 
				% % \item[]
			\end{compactdesc}
		\end{framed}

	% -----
	\subsubsection{Section 88:13} % *DC88-13 
	% -----
		\label{DC88-13}

		The light which is in all things, which giveth life to all things,%
			\footnote{%
				% \label{}%
				I interpret this part of the verse in the sense that we receive ``life'' from Christ in the way that Christ receives life from the Father; see \hyperref[JOHN-5-26]{John 5:26}, \hyperref[JOHN-6-56]{John 6:56--57}, sort of \hyperref[JOHN-10-18]{John 10:18}, \hyperref[DC93-2]{DC 93:2}, \hyperref[DC93-9]{DC 93:9}. There are probably other verses that are relevant to this as well.
				}
		which is the law by which all things are governed, even the power of God%
			\footnote{%
				% \label{}%
				To me it seems clear that the ``light'' being spoken of here and in the previous verse, and possibly the others, is the priesthood. It is the ``law by which all things are governed, even the power of God.'' This gives you a much more direct connection to BY-style ideas of the priesthood than you can find elsewhere.
				}
		who sitteth upon his throne, who is in the \hyperref[BOSOM]{bosom} of eternity,%
			\footnote{%
				% \label{}%
				This is the only occurrence of ``bosom of eternity'' in the scriptures; see \hyperref[BOSOM-PHRASES]{here}.
				} 
		who is in the midst of all things.

		\begin{framed}
		\scriptsize
			\begin{compactdesc}
				\item[JSP] the light which is in all things which giveth life to all things, which is the law by which all things are govorned, even the power of God, who sitteth upon his throne; who is in the bosom of eternity, who is in the midst of all things
				% % \item[BC] 
				% % \item[]
			\end{compactdesc}
		\end{framed}

	% -----
	\subsubsection{Section 88:14} % *DC88-14 
	% -----
		\label{DC88-14}

			\footnote{%
				% \label{}%
				We're clearly beginning a new division in this section at this verse. The idea of a ``redemption'' is somewhat implicit in verse 4, so that might be the connection to this new division.
				}%
		Now, verily I say unto you, that through the \hyperref[REDEMPTION]{redemption} which is made for you is brought to pass the \hyperref[RESURRECTION]{resurrection} from the dead.

		\begin{framed}
		\scriptsize
			\begin{compactdesc}
				\item[JSP] Now verily I say unto you, that through the redemption, which is made for you; is brought to pass the resurection from the dead;
				% % \item[BC] 
				% % \item[]
			\end{compactdesc}
		\end{framed}

	% -----
	\subsubsection{Section 88:15} % *DC88-15
	% -----
		\label{DC88-15}

		And the spirit and the body are the soul of man.

		\begin{framed}
		\scriptsize
			\begin{compactdesc}
				\item[JSP] (and the spirit, and the body is the soul of man)
				% % \item[BC] 
				% % \item[]
			\end{compactdesc}
		\end{framed}

	% -----
	\subsubsection{Section 88:16} % *DC88-16 
	% -----
		\label{DC88-16}

		And the \hyperref[RESURRECTION]{resurrection} from the dead is the \hyperref[REDEMPTION]{redemption} of the soul.

		\begin{framed}
		\scriptsize
			\begin{compactdesc}
				\item[JSP] and the resurection from the dead, is the redemption of the soul; and the redemption of the soul,
				% % \item[BC] 
				% % \item[]
			\end{compactdesc}
		\end{framed}

	% -----
	\subsubsection{Section 88:17} % *DC88-17 
	% -----
		\label{DC88-17}

		And the \hyperref[REDEMPTION]{redemption} of the soul is through him that \hyperref[QUICKEN]{quickeneth} all things,%
		\footnote{%
			% \label{}%
			See \hyperref[DC33-16]{DC 33:16}, ``and the power of my Spirit quickeneth all things.''
			} 
		in whose bosom it is decreed that the poor and the meek of the earth shall inherit it.

		\begin{framed}
		\scriptsize
			\begin{compactdesc}
				\item[JSP] is through him, who quickneth all things, in whose bosom, it is decreed, that the poor, and the meek of the earth, shall inherit it;
				% % \item[BC] 
				% % \item[]
			\end{compactdesc}
		\end{framed}

	% -----
	\subsubsection{Section 88:18} % *DC88-18 
	% -----
		\label{DC88-18}

		Therefore, it must needs be sanctified from all unrighteousness, that it may be prepared for the celestial glory;

		\begin{framed}
		\scriptsize
			\begin{compactdesc}
				\item[JSP] therefore it must needs be sanctified, from all unrighteousness, that it may be prepared for the celestial glory;
				% % \item[BC] 
				% % \item[]
			\end{compactdesc}
		\end{framed}

	% -----
	\subsubsection{Section 88:19} % *DC88-19 
	% -----
		\label{DC88-19}

		For after it hath filled the measure%
			\footnote{%
				% \label{}%
				Paul talks a lot about ``measures'' at \hyperref[2CORINTHIANS-10-13]{2 Corinthians 10:13--15}, though I'm not sure how related it is to this stuff.
				}
		of its creation,%
			\footnote{%
				\label{DC88-19-measure}%
				See \hyperref[DC88-25]{verse 25}, and the \hyperref[TEMPLE-ENDOWMENT-ENDOW-PRE-1990-DAY5]{fifth day} of the endowment ceremony. There's also a connection here to the idea of being ``made perfect'' as spoken of in \hyperref[DC76-69]{DC 76:69}---the source of some of the language there is \hyperref[HEBREWS-12-23]{Hebrews 12:23--24} and the \gr{τετελειωμένων} there, a participle from \hyperref[TELEIOO]{\grl{τελειόω}}. Definition \#2.b for \gr{τελειόω} is ``bring to full measure, fill the measure of,'' so perhaps we can read this verse here in DC 88 (and the other uses of the construction elsewhere) as ``being perfected'' or ``being made perfect'' in the way that Christ can do that.
				
				See also \hyperref[MATTHEW-23-32]{Matthew 23:32}; \hyperref[DC49-16]{DC 49:16--17}; \hyperref[DC103-3]{DC 103:3}; 
				}
		it shall be crowned with glory, even with the presence of God the Father;

		\begin{framed}
		\scriptsize
			\begin{compactdesc}
				\item[JSP] for after it hath filled the measure of its creation, it shall be crowned with \sout{the} glory, even with the presence of God the father;
				% % \item[BC] 
				% % \item[]
			\end{compactdesc}
		\end{framed}

	% -----
	\subsubsection{Section 88:20} % *DC88-20 
	% -----
		\label{DC88-20}

		That bodies%
			\footnote{%
				% \label{}%
				This ``bodies'' language is unusual, and reminds me of \hyperref[DC76-24]{DC 76:24}---maybe it's a way of speaking about any being who has passed through the requisite processes.
				}
		who are of the celestial kingdom may possess it forever and ever; for, for this intent was it made and created, and for this intent are they sanctified.%
			\footnote{%
				% \label{}%
				Compare \hyperref[1NEPHI-17-36]{1 Nephi 17:36}.
				}

		\begin{framed}
		\scriptsize
			\begin{compactdesc}
				\item[JSP] that bodies, who are of the celestial kingdom may posses it, for ever, \& ever; for, for this intent was it made, and created, and for this intent, are they sanctified,
				% % \item[BC] 
				% % \item[]
			\end{compactdesc}
		\end{framed}

	% -----
	\subsubsection{Section 88:21} % *DC88-21 
	% -----
		\label{DC88-21}

		And they who are not sanctified through the law which I have given unto you, even the law of Christ,%
			\footnote{%
				% \label{}%
				See more on the phrase ``law of Christ'' \hyperref[LAW-PHRASES]{here}. I think the phrase is a way of referring to the priesthood while calling attention to Christ's central role it (as \hyperref[DC107-3]{DC 107:3} suggests).
				}
		must inherit another kingdom, even that of a terrestrial kingdom, or that of a telestial kingdom.

		\begin{framed}
		\scriptsize
			\begin{compactdesc}
				\item[JSP] and they who are not sanctified, through the law which I have given unto you; even the law of Christ, must inherit another kingdom even that of a Terestrial kingdom, or that of a telestial kingdom,
				% % \item[BC] 
				% % \item[]
			\end{compactdesc}
		\end{framed}

	% -----
	\subsubsection{Section 88:22} % *DC88-22 
	% -----
		\label{DC88-22}

		For he who is not able to \hyperref[ABIDE]{abide}%
			\footnote{%
				% \label{}%
				The use of ``abide'' here could be ``abide by,'' just dropping the preposition; the 1828 has this to say about ``abide by'': ``In general, \textit{abide} by signifies to adhere to, maintain defend, or stand to, as to \textit{abide} by a promise, or by a friend; or to suffer the consequences, as to \textit{abide} by the event, that is, to be fixed or permanent in a particular condition.'' But this use could also be one of the other definitions of just ``abide,'' including \#3, ``To bear or endure; to bear patiently.'' In that case this verse would relate to \hyperref[DC132-16]{DC 132:16}. See \hyperref[DC88-26]{verse 26} below for a clear use of ``abide'' in this sense. See \hyperref[DC88-35]{verse 35} for a use of ``abide by,'' along with ``abide in''; ``abide'' is clearly a key concept in this section. See also \hyperref[DC132-6]{DC 132:6}; \hyperref[DC132-16]{DC 132:16--17}; 
				}
		the law of a celestial kingdom%
			\footnote{%
				% \label{}%
				We're talking about the ``celestial law'' here.
				
				See \hyperref[DC78-7]{DC 78:7} and \hyperref[DC78-13]{13}; 
				}
		cannot abide a celestial glory.

		\begin{framed}
		\scriptsize
			\begin{compactdesc}
				\item[JSP] for he that is not able to abide the law of a celestial kingdom cannot abide a celestial glory,
				% % \item[BC] 
				% % \item[]
			\end{compactdesc}
		\end{framed}

	% -----
	\subsubsection{Section 88:23} % *DC88-23 
	% -----
		\label{DC88-23}

		And he who cannot \hyperref[ABIDE]{abide} the law of a terrestrial kingdom cannot abide a terrestrial glory.

		\begin{framed}
		\scriptsize
			\begin{compactdesc}
				\item[JSP] and he who cannot abide, the law of a Terestrial kingdom cannot, abide a Terestrial glory,
				% % \item[BC] 
				% % \item[]
			\end{compactdesc}
		\end{framed}

	% -----
	\subsubsection{Section 88:24} % *DC88-24 
	% -----
		\label{DC88-24}

		And he who cannot \hyperref[ABIDE]{abide} the law of a telestial kingdom cannot \hyperref[ABIDE]{abide} a telestial glory; therefore he is not meet for a kingdom of glory. Therefore he must \hyperref[ABIDE]{abide} a kingdom which is not a kingdom of glory.

		\begin{framed}
		\scriptsize
			\begin{compactdesc}
				\item[JSP] he who cannot abide the law of a Telestial kingdom cannot abide a Telestial glory; therefore he is not meet, for a kingdom of Glory, therefore he must abide a kingdom, which is not a kingdom of glory.
				% % \item[BC] 
				% % \item[]
			\end{compactdesc}
		\end{framed}

	% -----
	\subsubsection{Section 88:25} % *DC88-25 
	% -----
		\label{DC88-25}

		And again, verily I say unto you, the earth \hyperref[ABIDE]{abideth} the law of a celestial kingdom, for it filleth the measure of its creation,%
			\footnote{%
				% \label{}%
				For ``filleth the measure of its creation,'' see \hyperref[DC88-19]{verse 19}.
				} 
		and transgresseth not the law---

		\begin{framed}
		\scriptsize
			\begin{compactdesc}
				\item[JSP] And again, verily I say unto you, the earth abideth the law, of a celestial kingdom, for it filleth, the measur of its creation; and transg[r]esseth not the law
				% % \item[BC] 
				% % \item[]
			\end{compactdesc}
		\end{framed}

	% -----
	\subsubsection{Section 88:26} % *DC88-26 
	% -----
		\label{DC88-26}

		Wherefore, it shall be sanctified; yea, \hyperref[notwithstanding]{notwithstanding} it shall die, it shall be \hyperref[QUICKEN]{quickened} again, and shall \hyperref[ABIDE]{abide} the power by which it is \hyperref[QUICKEN]{quickened}, and the righteous shall inherit it.

		\begin{framed}
		\scriptsize
			\begin{compactdesc}
				\item[JSP] wherefore it shall be sanctified, yea notwithstanding it shall die, it shall be quickened again, and shall abide the power, by which it was quickened, and the righteous shall inherit it,
				% % \item[BC] 
				% % \item[]
			\end{compactdesc}
		\end{framed}

	% -----
	\subsubsection{Section 88:27} % *DC88-27 
	% -----
		\label{DC88-27}

		For \hyperref[notwithstanding]{notwithstanding} they die, they also shall rise again, a spiritual body.%
			\footnote{%
				% \label{}%
				This use of ``spiritual'' is very important, because it tells us something about the nature of our bodies after the \hyperref[RESURRECTION]{resurrection}. See \hyperref[ALMA-11-45]{Alma 11:45}.
				}

		\begin{framed}
		\scriptsize
			\begin{compactdesc}
				\item[JSP] for notwithstanding they die, they also shall rise again, a spiritual body,
				% % \item[BC] 
				% % \item[]
			\end{compactdesc}
		\end{framed}

	% -----
	\subsubsection{Section 88:28} % *DC88-28 
	% -----
		\label{DC88-28}

		They who are of a celestial spirit shall receive the same body which was a natural body; even ye shall receive your bodies, and your glory shall be that glory by which your bodies are \hyperref[QUICKEN]{quickened}.%
			\footnote{%
				% \label{}%
				Verses 29--32 are examples and explanations of this last comment here, about what it means to receive the glory by which one's body is quickened.
				}

		\begin{framed}
		\scriptsize
			\begin{compactdesc}
				\item[JSP] they who are of a celestial spirit, shall receive the same body which was a natural body, even ye shall, receive your bodies; and your glory, shall be that glory, by which your bodies, are quickened,
				% % \item[BC] 
				% % \item[]
			\end{compactdesc}
		\end{framed}

	% -----
	\subsubsection{Section 88:29} % *DC88-29 
	% -----
		\label{DC88-29}

		Ye who are \hyperref[QUICKEN]{quickened} by a portion of the celestial glory shall then receive of the same, even a fulness.

		\begin{framed}
		\scriptsize
			\begin{compactdesc}
				\item[JSP] ye who are quickened, by a portion, of the celestial glory, shall then receive of the same, even a fulness,
				% % \item[BC] 
				% % \item[]
			\end{compactdesc}
		\end{framed}

	% -----
	\subsubsection{Section 88:30} % *DC88-30 
	% -----
		\label{DC88-30}

		And they who are \hyperref[QUICKEN]{quickened} by a portion of the terrestrial glory shall then receive of the same, even a fulness.

		\begin{framed}
		\scriptsize
			\begin{compactdesc}
				\item[JSP] and they, who are quickened, by a portion of the Terestriall glory, shall then, receive of the same even a fulness;
				% % \item[BC] 
				% % \item[]
			\end{compactdesc}
		\end{framed}

	% -----
	\subsubsection{Section 88:31} % *DC88-31 
	% -----
		\label{DC88-31}

		And also they who are \hyperref[QUICKEN]{quickened} by a portion of the telestial glory shall then receive of the same, even a fulness.

		\begin{framed}
		\scriptsize
			\begin{compactdesc}
				\item[JSP] and also they who are quickened by a portion, of the Telestial glory, shall then receive of the same, even a fulness,
				% % \item[BC] 
				% % \item[]
			\end{compactdesc}
		\end{framed}

	% -----
	\subsubsection{Section 88:32} % *DC88-32 
	% -----
		\label{DC88-32}

		And they who remain shall also be \hyperref[QUICKEN]{quickened}; nevertheless, they shall return again to their own place, to enjoy that which they are willing%
			\footnote{%
				% \label{}%
				Compare ``willing'' here with \hyperref[JOHN-5-35]{John 5:35}.
				}
		to receive, because they were not willing to enjoy that which they might have received.%
			\footnote{%
				% \label{DC88-32-NOTE-enjoy}%
				This is a really crucial idea---that you receive as much as you are willing to enjoy, or are willing to pay the price to enjoy. There are many, many reasons why someone might choose to receive less than they otherwise could:
					\begin{compactitem}
						\item They don't want what they have done or the ways they have lived to become known (\hyperref[JOHN-3-20]{John 3:20--21}).
					\end{compactitem}
				See also \hyperref[1NEPHI-17-21]{1 Nephi 17:21}, \hyperref[MOSIAH-26-25]{Mosiah 26:25--26}, \hyperref[MORMON-9-3]{Mormon 9:3--5}, \hyperref[MORMON-9-14]{Mormon 9:14}, 
				}

		\begin{framed}
		\scriptsize
			\begin{compactdesc}
				\item[JSP] and they who remain, shall also be quickend, nevertheless they shall, return again, to there own place, to enjoy that which they are willing to receive; because they were not, willing, to enjoy that which they might have received;
				% % \item[BC] 
				% % \item[]
			\end{compactdesc}
		\end{framed}

	% -----
	\subsubsection{Section 88:33} % *DC88-33 
	% -----
		\label{DC88-33}

		For what doth it profit a man if a gift%
			\footnote{%
				% \label{}%
				With what is coming next in verse 34, it could be that the ``gift'' spoken of here is the ``law of Christ,'' as discussed in \hyperref[DC88-21]{verse 21}.
				}
		is bestowed upon him, and he receive not the gift?%
			\footnote{%
				% \label{}%
				Compare \hyperref[MORONI-7-6]{Moroni 7:6--10}.
				}
		Behold, he rejoices not in that which is given unto him, neither rejoices in him who is the giver of the gift.

		\begin{framed}
		\scriptsize
			\begin{compactdesc}
				\item[JSP] for what doth it poffet [profit] a man, if a gift, is bestowed, upon him and he receive not the gift, behold he rejoiceth not, in that which is given unto him, neither rejoice in him, who is the giver of the gift;
				% % \item[BC] 
				% % \item[]
			\end{compactdesc}
		\end{framed}

	% -----
	\subsubsection{Section 88:34} % *DC88-34 
	% -----
		\label{DC88-34}

		And again, verily I say unto you, that which is governed by law%
			\footnote{%
				% \label{}%
				I read this ``law'' here as speaking about the whole structure of laws we call the priesthood.
				}
		is also preserved by law and perfected and sanctified by the same.%
			\footnote{%
				\label{DC88-34-law}%
				See \hyperref[LCHRIST-061]{event \#61} of Christ's life, where you get the words \gr{ἀνομία} (Matthew) and \gr{ἀδικία} (Luke); both are clearly connected to laws. Remember that \gr{δικαιοσύνη} is just \gr{δίκαιος}, ``just, lawful,'' with the -\gr{σύνη} suffix to make it an abstract noun. To be \gr{δίκαιος} is to engage in lawful conduct; to promote \gr{δικαιοσύνη} in your life is to live that same way. The preserving, perfecting, and sanctifying power is \hyperref[DUNAMIS]{\grl{δύναμις}}; compare \hyperref[DC84-20]{DC 84:20--22}.
				
				See also \hyperref[DC43-8]{DC 43:8--10}, \hyperref[DC44-4]{DC 44:4--6} (which includes the idea of being ``preserved'' in connection with laws), 
				}

		\begin{framed}
		\scriptsize
			\begin{compactdesc}
				\item[JSP] and again, verily I say unto you, that which is govorned by law, is also preserved by law, and perfected, and sanctified by the same;
				% % \item[BC] 
				% % \item[]
			\end{compactdesc}
		\end{framed}

	% -----
	\subsubsection{Section 88:35} % *DC88-35 
	% -----
		\label{DC88-35}

		That which breaketh a law, and abideth not by law,%
			\footnote{%
				% \label{}%
				See also \hyperref[ROMANS-8-7]{Romans 8:7}; \hyperref[DC98-8]{DC 98:8}; 
				}
		but seeketh to become a law unto%
			\footnote{%
				% \label{}%
				For ``law unto themselves,'' see \hyperref[ROMANS-2-14]{Romans 2:14}, \hyperref[ROMANS-10-3]{Romans 10:3}, 
				}
		itself,%
			\footnote{%
				% \label{}%
				I connect this ``become a law unto itself'' with the priesthood, specifically, with trying to become the person through whom the priesthood is received and governed in an implementation of the gospel. For us that is Christ, but Lucifer apparently wanted to be that individual---see \hyperref[MOSES-4-1]{Moses 4:1}, and connect to the idea of ``honor'' in \hyperref[HEBREWS-5-4]{Hebrews 5:4}.
				}
		and willeth to abide in sin, and altogether abideth in sin,
			\footnote{%
				% \label{}%
				Compare the list of things that cannot sanctify here to what Christ says in \hyperref[MATTHEW-23-23]{Matthew 23:23}.
				}%
		cannot be sanctified by law, neither by mercy, justice, nor judgment.%
			\footnote{%
				% \label{}%
				Compare \hyperref[HELAMAN-5-2]{Helaman 5:2--3}, 
				}
		Therefore, they must remain filthy still.

		\begin{framed}
		\scriptsize
			\begin{compactdesc}
				\item[JSP] that which breaketh a law, and abideth not by law, but seeketh to become a law unto itself, and willeth to abide in sin; and altogether abideth in sin, cannot be sanctified by law; neither of mercy, justice, or judgment; therefore they must remain filthy still,
				% % \item[BC] 
				% % \item[]
			\end{compactdesc}
		\end{framed}

	% -----
	\subsubsection{Section 88:36} % *DC88-36 
	% -----
		\label{DC88-36}

		All kingdoms have a law given;%
			\footnote{%
				% \label{}%
				In these and the next few verses, I want to give them a strong reading, like one that says that the laws of a kingdom constitute the kingdom---having a set of laws is what makes a kingdom a kingdom.
				} 

		\begin{framed}
		\scriptsize
			\begin{compactdesc}
				\item[JSP] all kingdoms have a law given;
				% % \item[BC] 
				% % \item[]
			\end{compactdesc}
		\end{framed}

	% -----
	\subsubsection{Section 88:37} % *DC88-37 
	% -----
		\label{DC88-37}

		And there are many kingdoms; for there is no space in the which there is no kingdom; and there is no kingdom in which there is no space, either a greater or a lesser kingdom.

		\begin{framed}
		\scriptsize
			\begin{compactdesc}
				\item[JSP] and there are many kingdoms; for there is no space, in the which there is no kingdom\sout{s}; and there is no kingdom, in which there is no space, eather a greater or lesser kingdom,
				% % \item[BC] 
				% % \item[]
			\end{compactdesc}
		\end{framed}

	% -----
	\subsubsection{Section 88:38} % *DC88-38 
	% -----
		\label{DC88-38}

		And unto every kingdom is given a law; and unto every law there are certain bounds%
			\footnote{%
				% \label{}%
				Compare \hyperref[DC121-30]{DC 121:30--31}.
				}
		also and conditions.%
			\footnote{%
				\label{DC88-38-NOTE-conditions}%
				For discussion of the idea of ``conditions'' for a law, see \hyperref[CONDITION-laws]{this study} at \hyperref[CONDITION]{Condition}. Think also of the ``\hyperref[CONDITION-PHRASES]{conditions of repentance}.''
				}

		\begin{framed}
		\scriptsize
			\begin{compactdesc}
				\item[JSP] and unto evry kingdom, is given a law, and unto evry law there are certain bounds also, and conditions,
				% % \item[BC] 
				% % \item[]
			\end{compactdesc}
		\end{framed}

	% -----
	\subsubsection{Section 88:39} % *DC88-39 
	% -----
		\label{DC88-39}

		All beings who abide not in those conditions are not justified.

		\begin{framed}
		\scriptsize
			\begin{compactdesc}
				\item[JSP] all beings who abide not, in those conditions, are not justified,
				% % \item[BC] 
				% % \item[]
			\end{compactdesc}
		\end{framed}

	% -----
	\subsubsection{Section 88:40} % *DC88-40 
	% -----
		\label{DC88-40}

			\footnote{%
				% \label{}%
				For a very similar idea to the one in this verse, see \hyperref[ALMA-41-12]{Alma 41:12--13}.
				}%
		For intelligence \hyperref[CLEAVE]{cleaveth}%
			\footnote{%
				% \label{}%
				Compare ``cleave'' in \hyperref[DC98-11]{DC 98:11}.
				}
		unto intelligence; wisdom receiveth wisdom; truth embraceth truth; virtue loveth virtue; light cleaveth unto light; mercy hath compassion on mercy and claimeth her own; justice continueth its course and claimeth its own; judgment goeth before the face of him who sitteth upon the throne and governeth and executeth all things.

		\begin{framed}
		\scriptsize
			\begin{compactdesc}
				\item[JSP] for inteligence, cleaveth unto inteligence, wisdom, receiveth wisdom, truth embraceth truth, virtue, Loveth virtue, light, Cleaveth unto light, mercy hath compassion on mercy, and claimeth her own; justice continueth its course and claimeth its own, judgment, goeth before the face of him, who sitteth up[on] the throne, and gove[r]neth, and executeth all things,
				% % \item[BC] 
				% % \item[]
			\end{compactdesc}
		\end{framed}

	% -----
	\subsubsection{Section 88:41} % *DC88-41 
	% -----
		\label{DC88-41}

		He comprehendeth all things, and all things are before him, and all things are round about him; and he is above all things, and in all things, and is through all things, and is round about all things; and all things are by him, and of him, even God, forever and ever.

		\begin{framed}
		\scriptsize
			\begin{compactdesc}
				\item[JSP] he comprehendeth all things, and all things, are before him, and all things, are round about him, and he is, above all things, and in all things, and is through all things, and is round about all thing[s] and all things are by him, and of him, even God for ever, and ever,
				% % \item[BC] 
				% % \item[]
			\end{compactdesc}
		\end{framed}

	% -----
	\subsubsection{Section 88:42} % *DC88-42 
	% -----
		\label{DC88-42}

		And again, verily I say unto you, he hath given a law%
			\footnote{%
				% \label{}%
				The priesthood. See \hyperref[PRIESTHOOD-BY]{BY on the priesthood}; probably the clearest one is from \hyperref[EMS-1854-BY-19-NOV-1854]{19 November 1854}: ``Now let me tell you there is no world made without that world has motion, and it is governed and controled by the law of the eternal preisthood; and if you want to know what preisthood is, it is a perfect system of Government by which planets and all creations are governed.''
				}
		unto all things, by which they move in their times and their seasons;

		\begin{framed}
		\scriptsize
			\begin{compactdesc}
				\item[JSP] And again <verily> I say unto you, he hath given a law, unto all things, by which they moove in there times, and there seasons,
				% % \item[BC] 
				% % \item[]
			\end{compactdesc}
		\end{framed}

	% -----
	\subsubsection{Section 88:43} % *DC88-43 
	% -----
		\label{DC88-43}

		And their courses are fixed, even the courses of the heavens and the earth, which comprehend the earth and all the planets.%
			\footnote{%
				% \label{}%
				Why leave out the stars and everything else?
				}

		\begin{framed}
		\scriptsize
			\begin{compactdesc}
				\item[JSP] and there cources are fixed, even the cources of the heavens, and the earth which, comprehend the earth, and all the planets,
				% % \item[BC] 
				% % \item[]
			\end{compactdesc}
		\end{framed}

	% -----
	\subsubsection{Section 88:44} % *DC88-44 
	% -----
		\label{DC88-44}

		And they give light to each other in their times and in their seasons, in their minutes, in their hours, in their days, in their weeks, in their months, in their years---all these are one year with God, but not with man.

		\begin{framed}
		\scriptsize
			\begin{compactdesc}
				\item[JSP] and they give light to each other in there times, and in there seasons, in there minuits in there hours, in there days, in there week, in there months, in there years; all these are one year with God. but not with man;
				% % \item[BC] 
				% % \item[]
			\end{compactdesc}
		\end{framed}

	% -----
	\subsubsection{Section 88:45} % *DC88-45 
	% -----
		\label{DC88-45}

		The earth rolls%
			\footnote{%
				% \label{}%
				Compare \hyperref[DC121-30]{DC 121:30--31}.
				}
		upon her \hyperref[WING]{wings},%
			\footnote{%
				% \label{}%
				A very unusual reference to ``wings'' here; I'm not sure what it means. Check the 1828 definitions \hyperref[WING]{here}; I'm thinking it's definition \#4, \#5, or \#6. \#6 is ``Motive or incitement of flight,'' so it could be referring to the law from \hyperref[DC88-42]{verse 42}.
				}
		and the sun giveth his light by day, and the moon giveth her light by night, and the stars also give their light, as they roll upon their \hyperref[WING]{wings} in their glory, in the midst of the power of God.

		\begin{framed}
		\scriptsize
			\begin{compactdesc}
				\item[JSP] the <Earth> rolls upon her wings, and the sun giveth her light by day, and the moon, giveth her light by night, and the stars also giveth there light as they roll upon, there wings, in there glory in the midst, of the power, of God,
				% % \item[BC] 
				% % \item[]
			\end{compactdesc}
		\end{framed}

	% -----
	\subsubsection{Section 88:46} % *DC88-46 
	% -----
		\label{DC88-46}

		Unto what shall I liken these kingdoms, that ye may understand?

		\begin{framed}
		\scriptsize
			\begin{compactdesc}
				\item[JSP] unto what shall I liken these kingdoms, that ye may understand,
				% % \item[BC] 
				% % \item[]
			\end{compactdesc}
		\end{framed}

	% -----
	\subsubsection{Section 88:47} % *DC88-47 
	% -----
		\label{DC88-47}

		Behold, all these are kingdoms, and any man who hath seen any or the least of these hath seen God moving in his majesty and power.

		\begin{framed}
		\scriptsize
			\begin{compactdesc}
				\item[JSP] behold all these <are> kingdoms, and any man, who hath seen, any, or the least of these, have seen God, moving in his magesty and power;
				% % \item[BC] 
				% % \item[]
			\end{compactdesc}
		\end{framed}

	% -----
	\subsubsection{Section 88:48} % *DC88-48 
	% -----
		\label{DC88-48}

		I say unto you, he hath seen him; nevertheless, he who came unto his own was not comprehended.

		\begin{framed}
		\scriptsize
			\begin{compactdesc}
				\item[JSP] I say unto you, he hath seen him, nevertheless, he who came unto his own, was not comprehended,
				% % \item[BC] 
				% % \item[]
			\end{compactdesc}
		\end{framed}

	% -----
	\subsubsection{Section 88:49} % *DC88-49 
	% -----
		\label{DC88-49}

		The light shineth in darkness, and the darkness comprehendeth it not; nevertheless, the day shall come when you shall comprehend even God, being \hyperref[QUICKEN]{quickened} in him and by him.

		\begin{framed}
		\scriptsize
			\begin{compactdesc}
				\item[JSP] the light shineth in darkness, and the darkness compre[hen]deth it not, nevertheless, the day shall come, when you shall, comprehend even God, being quickened in him, and by him,
				% % \item[BC] 
				% % \item[]
			\end{compactdesc}
		\end{framed}

	% -----
	\subsubsection{Section 88:50} % *DC88-50 
	% -----
		\label{DC88-50}

		Then shall ye know that ye have seen me, that I am, and that I am the true light that is in you, and that you are in me; otherwise ye could not \hyperref[ABOUND]{abound}.

		\begin{framed}
		\scriptsize
			\begin{compactdesc}
				\item[JSP] then shall ye know, that ye have seen me, that I am, and that I am the true light, that is in you, and that you are in me, otherwise ye could not abound,
				% % \item[BC] 
				% % \item[]
			\end{compactdesc}
		\end{framed}

	% -----
	\subsubsection{Section 88:51} % *DC88-51 
	% -----
		\label{DC88-51}

			\footnote{%
				% \label{}%
				The parable that the Lord is about to give is a way of expressing how God gives light to different kingdoms, or different organizations; again, I think this light is the priesthood and the laws which constitute it. \hyperref[DC88-48]{Verse 58}, for example, is calling back to verses 38, 44, and especially 12--13. Notice the reference also to ``order'' in \hyperref[DC88-60]{verse 60}. I think you can take the different kingdoms as either all being on Earth, or being on Earth and on other planets. If you go the former way, then you can see each different ``kingdom'' as a different time period in which different people were alive and the Lord came to manifest his light in different ways, ``every man in his own order.'' Notice that the Lord calls those receiving the revelation ``the first laborers in this last kingdom'' (\hyperref[DC88-70]{verse 70}, referring back to this parable (see also 74). I think you could also take the ``kingdoms'' as being on different planets, and still run a similar interpretation.
				}%
		Behold, I will liken these kingdoms unto a man having a field, and he sent forth his servants into the field to dig in the field.

		% \begin{framed}
		% \scriptsize
			% \begin{compactdesc}
				% \item[JSP] 
				% % \item[BC] 
				% % \item[]
			% \end{compactdesc}
		% \end{framed}

	% -----
	\subsubsection{Section 88:52} % *DC88-52 
	% -----
		\label{DC88-52}

		And he said unto the first: Go ye and labor in the field, and in the first hour I will come unto you, and ye shall behold the joy of my countenance.

		% \begin{framed}
		% \scriptsize
			% \begin{compactdesc}
				% \item[JSP] 
				% % \item[BC] 
				% % \item[]
			% \end{compactdesc}
		% \end{framed}

	% -----
	\subsubsection{Section 88:53} % *DC88-53 
	% -----
		\label{DC88-53}

		And he said unto the second: Go ye also into the field, and in the second hour I will visit you with the joy of my countenance.

		% \begin{framed}
		% \scriptsize
			% \begin{compactdesc}
				% \item[JSP] 
				% % \item[BC] 
				% % \item[]
			% \end{compactdesc}
		% \end{framed}

	% -----
	\subsubsection{Section 88:54} % *DC88-54 
	% -----
		\label{DC88-54}

		And also unto the third, saying: I will visit you;

		% \begin{framed}
		% \scriptsize
			% \begin{compactdesc}
				% \item[JSP] 
				% % \item[BC] 
				% % \item[]
			% \end{compactdesc}
		% \end{framed}

	% -----
	\subsubsection{Section 88:55} % *DC88-55 
	% -----
		\label{DC88-55}

		And unto the fourth, and so on unto the twelfth.

		% \begin{framed}
		% \scriptsize
			% \begin{compactdesc}
				% \item[JSP] 
				% % \item[BC] 
				% % \item[]
			% \end{compactdesc}
		% \end{framed}

	% -----
	\subsubsection{Section 88:56} % *DC88-56 
	% -----
		\label{DC88-56}

		And the lord of the field went unto the first in the first hour, and tarried with him all that hour, and he was made glad with the light of the countenance of his lord.

		% \begin{framed}
		% \scriptsize
			% \begin{compactdesc}
				% \item[JSP] 
				% % \item[BC] 
				% % \item[]
			% \end{compactdesc}
		% \end{framed}

	% -----
	\subsubsection{Section 88:57} % *DC88-57 
	% -----
		\label{DC88-57}

		And then he withdrew from the first that he might visit the second also, and the third, and the fourth, and so on unto the twelfth.

		% \begin{framed}
		% \scriptsize
			% \begin{compactdesc}
				% \item[JSP] 
				% % \item[BC] 
				% % \item[]
			% \end{compactdesc}
		% \end{framed}

	% -----
	\subsubsection{Section 88:58} % *DC88-58 
	% -----
		\label{DC88-58}

		And thus they all received the light of the countenance of their lord, every man in his hour, and in his time, and in his season---

		% \begin{framed}
		% \scriptsize
			% \begin{compactdesc}
				% \item[JSP] 
				% % \item[BC] 
				% % \item[]
			% \end{compactdesc}
		% \end{framed}

	% -----
	\subsubsection{Section 88:59} % *DC88-59 
	% -----
		\label{DC88-59}

		Beginning at the first, and so on unto the last, and from the last unto the first, and from the first unto the last;

		% \begin{framed}
		% \scriptsize
			% \begin{compactdesc}
				% \item[JSP] 
				% % \item[BC] 
				% % \item[]
			% \end{compactdesc}
		% \end{framed}

	% -----
	\subsubsection{Section 88:60} % *DC88-60 
	% -----
		\label{DC88-60}

		Every man in his own \hyperref[ORDER]{order}, until his hour was finished, even according as his lord had commanded him, that his lord might be glorified in him, and he in his lord, that they all might be glorified.

		% \begin{framed}
		% \scriptsize
			% \begin{compactdesc}
				% \item[JSP] 
				% % \item[BC] 
				% % \item[]
			% \end{compactdesc}
		% \end{framed}

	% -----
	\subsubsection{Section 88:61} % *DC88-61 
	% -----
		\label{DC88-61}

		Therefore, unto this parable I will liken all these kingdoms, and the inhabitants thereof---every kingdom in its hour, and in its time, and in its season, even according to the decree which God hath made.

		% \begin{framed}
		% \scriptsize
			% \begin{compactdesc}
				% \item[JSP] 
				% % \item[BC] 
				% % \item[]
			% \end{compactdesc}
		% \end{framed}

	% -----
	\subsubsection{Section 88:62} % *DC88-62 
	% -----
		\label{DC88-62}

		And again, verily I say unto you, my friends, I leave these sayings with you to ponder in your hearts, with this commandment which I give unto you, that ye shall call upon me while I am near---

		% \begin{framed}
		% \scriptsize
			% \begin{compactdesc}
				% \item[JSP] 
				% % \item[BC] 
				% % \item[]
			% \end{compactdesc}
		% \end{framed}

	% -----
	\subsubsection{Section 88:63} % *DC88-63 
	% -----
		\label{DC88-63}

		Draw near unto me and I will draw near unto you;%
			\footnote{%
				\label{DC88-63-NOTE-draw}%
				See my note at \hyperref[MOSIAH-2-36]{Mosiah 2:36}. I think the right way to think about this verse is in terms of Christ always being there and willing to meet us, but I'm not sure it's right to think of him as withdrawing in our less obedient phases. I think the right image is the one I describe in that note.
				}
		seek me diligently and ye shall find me; ask, and ye shall receive; knock, and it shall be opened unto you.

		% \begin{framed}
		% \scriptsize
			% \begin{compactdesc}
				% \item[JSP] 
				% % \item[BC] 
				% % \item[]
			% \end{compactdesc}
		% \end{framed}

	% -----
	\subsubsection{Section 88:64} % *DC88-64 
	% -----
		\label{DC88-64}

		Whatsoever ye ask the Father in my name it shall be given unto you, that is expedient for you;

		% \begin{framed}
		% \scriptsize
			% \begin{compactdesc}
				% \item[JSP] 
				% % \item[BC] 
				% % \item[]
			% \end{compactdesc}
		% \end{framed}

	% -----
	\subsubsection{Section 88:65} % *DC88-65 
	% -----
		\label{DC88-65}

		And if ye ask anything that is not expedient for you, it shall turn unto your condemnation.

		% \begin{framed}
		% \scriptsize
			% \begin{compactdesc}
				% \item[JSP] 
				% % \item[BC] 
				% % \item[]
			% \end{compactdesc}
		% \end{framed}

	% -----
	\subsubsection{Section 88:66} % *DC88-66 
	% -----
		\label{DC88-66}

		Behold, that which you hear is as the voice of one crying in the wilderness---in the wilderness, because you cannot see him---my voice, because my voice is Spirit; my Spirit is truth; truth \hyperref[ABIDE]{abideth} and hath no end; and if it be in you it shall \hyperref[ABOUND]{abound}.

		% \begin{framed}
		% \scriptsize
			% \begin{compactdesc}
				% \item[JSP] 
				% % \item[BC] 
				% % \item[]
			% \end{compactdesc}
		% \end{framed}

	% -----
	\subsubsection{Section 88:67} % *DC88-67 
	% -----
		\label{DC88-67}

		And if your eye be \hyperref[SINGLE]{single} to my glory, your whole bodies shall be filled with light, and there shall be no darkness in you; and that body which is filled with light comprehendeth all things.

		% \begin{framed}
		% \scriptsize
			% \begin{compactdesc}
				% \item[JSP] 
				% % \item[BC] 
				% % \item[]
			% \end{compactdesc}
		% \end{framed}

	% -----
	\subsubsection{Section 88:68} % *DC88-68 
	% -----
		\label{DC88-68}

		Therefore, \hyperref[SANCTIFICATION]{sanctify} yourselves that your minds become \hyperref[SINGLE]{single} to God, and the days will come that you shall see him; for he will unveil his face unto you, and it shall be in his own time, and in his own way, and according to his own will.

		% \begin{framed}
		% \scriptsize
			% \begin{compactdesc}
				% \item[JSP] 
				% % \item[BC] 
				% % \item[]
			% \end{compactdesc}
		% \end{framed}

	% -----
	\subsubsection{Section 88:69} % *DC88-69 
	% -----
		\label{DC88-69}

		Remember the great and last promise which I have made unto you; cast away your idle thoughts and your excess of laughter far from you.

		% \begin{framed}
		% \scriptsize
			% \begin{compactdesc}
				% \item[JSP] 
				% % \item[BC] 
				% % \item[]
			% \end{compactdesc}
		% \end{framed}

	% -----
	\subsubsection{Section 88:70} % *DC88-70 
	% -----
		\label{DC88-70}

		Tarry ye, tarry ye in this place, and call a solemn assembly, even of those who are the first laborers in this last kingdom.

		% \begin{framed}
		% \scriptsize
			% \begin{compactdesc}
				% \item[JSP] 
				% % \item[BC] 
				% % \item[]
			% \end{compactdesc}
		% \end{framed}

	% -----
	\subsubsection{Section 88:71} % *DC88-71 
	% -----
		\label{DC88-71}

		And let those whom they have warned in their traveling call on the Lord, and ponder the warning in their hearts which they have received, for a little season.

		% \begin{framed}
		% \scriptsize
			% \begin{compactdesc}
				% \item[JSP] 
				% % \item[BC] 
				% % \item[]
			% \end{compactdesc}
		% \end{framed}

	% -----
	\subsubsection{Section 88:72} % *DC88-72 
	% -----
		\label{DC88-72}

		Behold, and lo, I will take care of your flocks, and will raise up elders and send unto them.

		% \begin{framed}
		% \scriptsize
			% \begin{compactdesc}
				% \item[JSP] 
				% % \item[BC] 
				% % \item[]
			% \end{compactdesc}
		% \end{framed}

	% -----
	\subsubsection{Section 88:73} % *DC88-73 
	% -----
		\label{DC88-73}

		Behold, I will hasten my work in its time.

		% \begin{framed}
		% \scriptsize
			% \begin{compactdesc}
				% \item[JSP] 
				% % \item[BC] 
				% % \item[]
			% \end{compactdesc}
		% \end{framed}

	% -----
	\subsubsection{Section 88:74} % *DC88-74 
	% -----
		\label{DC88-74}

		And I give unto you, who are the first laborers in this last kingdom, a commandment that you assemble yourselves together, and organize yourselves, and prepare yourselves, and sanctify yourselves; yea, purify your hearts, and cleanse your hands and your feet before me, that I may make you clean;

		% \begin{framed}
		% \scriptsize
			% \begin{compactdesc}
				% \item[JSP] 
				% % \item[BC] 
				% % \item[]
			% \end{compactdesc}
		% \end{framed}

	% -----
	\subsubsection{Section 88:75} % *DC88-75 
	% -----
		\label{DC88-75}

		That I may testify unto your Father, and your God, and my God, that you are clean from the blood of this wicked generation;%
			\footnote{%
				% \label{}%
				This verse and the previous are similar to the \hyperref[TEMPLE-ENDOWMENT-ENDOW-PRE-1990-INTRO]{introduction} to the endowment. In fact, I've noticed a lot of language similar to that found in the endowment ceremony; I think the endowment is drawing from a lot of what is said in this section. See also \hyperref[DC88-85]{verse 85} and \hyperref[DC88-138]{verses 138--140}, which connect to the \hyperref[TEMPLE-ENDOWMENT-INIT-PRE-2005]{initiatory}.
				}
		that I may fulfil this promise, this great and last promise, which I have made unto you, when I will.

		% \begin{framed}
		% \scriptsize
			% \begin{compactdesc}
				% \item[JSP] 
				% % \item[BC] 
				% % \item[]
			% \end{compactdesc}
		% \end{framed}

	% -----
	\subsubsection{Section 88:76} % *DC88-76 
	% -----
		\label{DC88-76}

		Also, I give unto you a commandment that ye shall continue in prayer and fasting from this time forth.

		% \begin{framed}
		% \scriptsize
			% \begin{compactdesc}
				% \item[JSP] 
				% % \item[BC] 
				% % \item[]
			% \end{compactdesc}
		% \end{framed}

	% -----
	\subsubsection{Section 88:77} % *DC88-77 
	% -----
		\label{DC88-77}

		And I give unto you a commandment that you shall teach one another the doctrine of the kingdom.

		% \begin{framed}
		% \scriptsize
			% \begin{compactdesc}
				% \item[JSP] 
				% % \item[BC] 
				% % \item[]
			% \end{compactdesc}
		% \end{framed}

	% -----
	\subsubsection{Section 88:78} % *DC88-78 
	% -----
		\label{DC88-78}

		Teach ye diligently and my grace shall attend you, that you may be instructed more perfectly in theory, in principle, in doctrine, in the law of the gospel,%
			\footnote{%
				% \label{}%
				This is one of the laws we receive in the \hyperref[TEMPLE-ENDOWMENT-ENDOW-PRE-1990-GOSPEL]{endowment} ceremony. Note that that portion of the endowment then gives instruction about behavior, some of which seems to draw from here in DC 88.
				}
		in all things that pertain unto the kingdom of God, that are expedient for you to understand;

		% \begin{framed}
		% \scriptsize
			% \begin{compactdesc}
				% \item[JSP] 
				% % \item[BC] 
				% % \item[]
			% \end{compactdesc}
		% \end{framed}

	% -----
	\subsubsection{Section 88:79} % *DC88-79 
	% -----
		\label{DC88-79}

		Of things both in heaven and in the earth, and under the earth; things which have been, things which are, things which must shortly come to pass; things which are at home, things which are abroad; the wars and the perplexities of the nations, and the judgments which are on the land; and a knowledge also of countries and of kingdoms---

		% \begin{framed}
		% \scriptsize
			% \begin{compactdesc}
				% \item[JSP] 
				% % \item[BC] 
				% % \item[]
			% \end{compactdesc}
		% \end{framed}

	% -----
	\subsubsection{Section 88:80} % *DC88-80 
	% -----
		\label{DC88-80}

		That ye may be prepared in all things when I shall send you again to \hyperref[MAGNIFY]{magnify} the calling whereunto I have called you, and the mission with which I have commissioned you.

		% \begin{framed}
		% \scriptsize
			% \begin{compactdesc}
				% \item[JSP] 
				% % \item[BC] 
				% % \item[]
			% \end{compactdesc}
		% \end{framed}

	% -----
	\subsubsection{Section 88:81} % *DC88-81 
	% -----
		\label{DC88-81}

		Behold, I sent you out to testify and warn the people, and it becometh every man who hath been warned to warn his neighbor.

		% \begin{framed}
		% \scriptsize
			% \begin{compactdesc}
				% \item[JSP] 
				% % \item[BC] 
				% % \item[]
			% \end{compactdesc}
		% \end{framed}

	% -----
	\subsubsection{Section 88:82} % *DC88-82 
	% -----
		\label{DC88-82}

		Therefore, they are left without excuse, and their sins are upon their own heads.

		% \begin{framed}
		% \scriptsize
			% \begin{compactdesc}
				% \item[JSP] 
				% % \item[BC] 
				% % \item[]
			% \end{compactdesc}
		% \end{framed}

	% -----
	\subsubsection{Section 88:83} % *DC88-83 
	% -----
		\label{DC88-83}

		He that seeketh me early%
			\footnote{%
				% \label{}%
				Not sure what the meaning of this word is here. Could be referring to the parable from earlier?
				}
		shall find me, and shall not be forsaken.

		% \begin{framed}
		% \scriptsize
			% \begin{compactdesc}
				% \item[JSP] 
				% % \item[BC] 
				% % \item[]
			% \end{compactdesc}
		% \end{framed}

	% -----
	\subsubsection{Section 88:84} % *DC88-84 
	% -----
		\label{DC88-84}

		Therefore, tarry ye, and labor diligently, that you may be perfected in your ministry to go forth among the Gentiles for the last time, as many as the mouth of the Lord shall name, to bind up the law and seal up the testimony,%
			\footnote{%
				% \label{}%
				This is a weird set of phrases, since it's the exact reverse of what we read in \hyperref[ISAIAH-8-16]{Isaiah 8:16}.
				}
		and to prepare the saints for the hour of judgment which is to come;

		% \begin{framed}
		% \scriptsize
			% \begin{compactdesc}
				% \item[JSP] 
				% % \item[BC] 
				% % \item[]
			% \end{compactdesc}
		% \end{framed}

	% -----
	\subsubsection{Section 88:85} % *DC88-85 
	% -----
		\label{DC88-85}

		That their souls may escape the wrath of God, the desolation of abomination which awaits the wicked, both in this world and in the world to come. Verily, I say unto you, let those who are not the first elders continue in the vineyard until the mouth of the Lord shall call them, for their time is not yet come; their garments are not clean from the blood of this generation.

		% \begin{framed}
		% \scriptsize
			% \begin{compactdesc}
				% \item[JSP] 
				% % \item[BC] 
				% % \item[]
			% \end{compactdesc}
		% \end{framed}

	% -----
	\subsubsection{Section 88:86} % *DC88-86 
	% -----
		\label{DC88-86}

		Abide ye in the liberty%
			\footnote{%
				% \label{}%
				Compare \hyperref[2NEPHI-2-27]{2 Nephi 2:27}.
				} 
		wherewith ye are made free; entangle not yourselves in sin, but let your hands be clean, until the Lord comes.

		% \begin{framed}
		% \scriptsize
			% \begin{compactdesc}
				% \item[JSP] 
				% % \item[BC] 
				% % \item[]
			% \end{compactdesc}
		% \end{framed}

	% -----
	\subsubsection{Section 88:87} % *DC88-87 
	% -----
		\label{DC88-87}

		For not many days hence and the earth shall tremble and reel to and fro as a drunken man; and the sun shall hide his face, and shall refuse to give light; and the moon shall be bathed in blood; and the stars shall become exceedingly angry, and shall cast themselves down as a fig that falleth from off a fig-tree.

		% \begin{framed}
		% \scriptsize
			% \begin{compactdesc}
				% \item[JSP] 
				% % \item[BC] 
				% % \item[]
			% \end{compactdesc}
		% \end{framed}

	% -----
	\subsubsection{Section 88:88} % *DC88-88 
	% -----
		\label{DC88-88}

		And after your testimony cometh wrath and indignation upon the people.

		% \begin{framed}
		% \scriptsize
			% \begin{compactdesc}
				% \item[JSP] 
				% % \item[BC] 
				% % \item[]
			% \end{compactdesc}
		% \end{framed}

	% -----
	\subsubsection{Section 88:89} % *DC88-89 
	% -----
		\label{DC88-89}

		For after your testimony cometh the testimony of earthquakes, that shall cause groanings in the midst of her, and men shall fall upon the ground and shall not be able to stand.

		% \begin{framed}
		% \scriptsize
			% \begin{compactdesc}
				% \item[JSP] 
				% % \item[BC] 
				% % \item[]
			% \end{compactdesc}
		% \end{framed}

	% -----
	\subsubsection{Section 88:90} % *DC88-90 
	% -----
		\label{DC88-90}

		And also cometh the testimony of the voice of thunderings, and the voice of lightnings, and the voice of tempests, and the voice of the waves of the sea heaving themselves beyond their bounds.

		% \begin{framed}
		% \scriptsize
			% \begin{compactdesc}
				% \item[JSP] 
				% % \item[BC] 
				% % \item[]
			% \end{compactdesc}
		% \end{framed}

	% -----
	\subsubsection{Section 88:91} % *DC88-91 
	% -----
		\label{DC88-91}

		And all things shall be in commotion; and surely, men's hearts shall fail them; for fear shall come upon all people.

		% \begin{framed}
		% \scriptsize
			% \begin{compactdesc}
				% \item[JSP] 
				% % \item[BC] 
				% % \item[]
			% \end{compactdesc}
		% \end{framed}

	% -----
	\subsubsection{Section 88:92} % *DC88-92 
	% -----
		\label{DC88-92}

		And angels shall fly through the midst of heaven, crying with a loud voice, sounding the trump of God, saying: Prepare ye, prepare ye, O inhabitants of the earth; for the judgment of our God is come. Behold, and lo, the Bridegroom cometh; go ye out to meet him.

		% \begin{framed}
		% \scriptsize
			% \begin{compactdesc}
				% \item[JSP] 
				% % \item[BC] 
				% % \item[]
			% \end{compactdesc}
		% \end{framed}

	% -----
	\subsubsection{Section 88:93} % *DC88-93 
	% -----
		\label{DC88-93}

		And immediately there shall appear a great sign in heaven, and all people shall see it together.

		% \begin{framed}
		% \scriptsize
			% \begin{compactdesc}
				% \item[JSP] 
				% % \item[BC] 
				% % \item[]
			% \end{compactdesc}
		% \end{framed}

	% -----
	\subsubsection{Section 88:94} % *DC88-94 
	% -----
		\label{DC88-94}

		And another angel shall sound his trump, saying: That great church, the mother of abominations, that made all nations drink of the wine of the wrath of her fornication, that persecuteth the saints of God, that shed their blood---she who sitteth upon many waters, and upon the islands of the sea---behold, she is the tares of the earth; she is bound in bundles; her bands are made strong, no man can loose them; therefore, she is ready to be burned. And he shall sound his trump both long and loud, and all nations shall hear it.

		% \begin{framed}
		% \scriptsize
			% \begin{compactdesc}
				% \item[JSP] 
				% % \item[BC] 
				% % \item[]
			% \end{compactdesc}
		% \end{framed}

	% -----
	\subsubsection{Section 88:95} % *DC88-95 
	% -----
		\label{DC88-95}

		And there shall be silence in heaven for the space of half an hour; and immediately after shall the curtain of heaven be unfolded, as a scroll is unfolded after it is rolled up, and the face of the Lord shall be unveiled;

		% \begin{framed}
		% \scriptsize
			% \begin{compactdesc}
				% \item[JSP] 
				% % \item[BC] 
				% % \item[]
			% \end{compactdesc}
		% \end{framed}

	% -----
	\subsubsection{Section 88:96} % *DC88-96 
	% -----
		\label{DC88-96}

		And the saints that are upon the earth, who are alive, shall be \hyperref[QUICKEN]{quickened} and be caught up to meet him.

		% \begin{framed}
		% \scriptsize
			% \begin{compactdesc}
				% \item[JSP] 
				% % \item[BC] 
				% % \item[]
			% \end{compactdesc}
		% \end{framed}

	% -----
	\subsubsection{Section 88:97} % *DC88-97 
	% -----
		\label{DC88-97}

		And they who have slept in their graves shall come forth, for their graves shall be opened; and they also shall be caught up to meet him in the midst of the pillar of heaven---

		% \begin{framed}
		% \scriptsize
			% \begin{compactdesc}
				% \item[JSP] 
				% % \item[BC] 
				% % \item[]
			% \end{compactdesc}
		% \end{framed}

	% -----
	\subsubsection{Section 88:98} % *DC88-98 
	% -----
		\label{DC88-98}

		They are Christ's, the \hyperref[FIRSTFRUITS]{first fruits}, they who shall descend with him first, and they who are on the earth and in their graves, who are first caught up to meet him; and all this by the voice of the sounding of the trump of the angel of God.

		% \begin{framed}
		% \scriptsize
			% \begin{compactdesc}
				% \item[JSP] 
				% % \item[BC] 
				% % \item[]
			% \end{compactdesc}
		% \end{framed}

	% -----
	\subsubsection{Section 88:99} % *DC88-99 
	% -----
		\label{DC88-99}

		And after this another angel shall sound, which is the second trump; and then cometh the redemption of those who are Christ's at his coming; who have received their part in that prison which is prepared for them, that they might receive the gospel, and be judged according to men in the flesh.

		% \begin{framed}
		% \scriptsize
			% \begin{compactdesc}
				% \item[JSP] 
				% % \item[BC] 
				% % \item[]
			% \end{compactdesc}
		% \end{framed}

	% -----
	\subsubsection{Section 88:100} % *DC88-100 
	% -----
		\label{DC88-100}

		And again, another trump shall sound, which is the third trump; and then come the spirits of men who are to be judged, and are found under condemnation;

		% \begin{framed}
		% \scriptsize
			% \begin{compactdesc}
				% \item[JSP] 
				% % \item[BC] 
				% % \item[]
			% \end{compactdesc}
		% \end{framed}

	% -----
	\subsubsection{Section 88:101} % *DC88-101 
	% -----
		\label{DC88-101}

		And these are the rest of the dead; and they live not again until the thousand years are ended, neither again, until the end of the earth.

		% \begin{framed}
		% \scriptsize
			% \begin{compactdesc}
				% \item[JSP] 
				% % \item[BC] 
				% % \item[]
			% \end{compactdesc}
		% \end{framed}

	% -----
	\subsubsection{Section 88:102} % *DC88-102 
	% -----
		\label{DC88-102}

		And another trump shall sound, which is the fourth trump, saying: There are found among those who are to remain until that great and last day, even the end, who shall remain filthy still.

		% \begin{framed}
		% \scriptsize
			% \begin{compactdesc}
				% \item[JSP] 
				% % \item[BC] 
				% % \item[]
			% \end{compactdesc}
		% \end{framed}

	% -----
	\subsubsection{Section 88:103} % *DC88-103 
	% -----
		\label{DC88-103}

		And another trump shall sound, which is the fifth trump, which is the fifth angel who committeth the everlasting gospel---flying through the midst of heaven, unto all nations, kindreds, tongues, and people;

		% \begin{framed}
		% \scriptsize
			% \begin{compactdesc}
				% \item[JSP] 
				% % \item[BC] 
				% % \item[]
			% \end{compactdesc}
		% \end{framed}

	% -----
	\subsubsection{Section 88:104} % *DC88-104 
	% -----
		\label{DC88-104}

		And this shall be the sound of his trump, saying to all people, both in heaven and in earth, and that are under the earth---for every ear shall hear it, and every knee shall bow, and every tongue shall confess, while they hear the sound of the trump, saying: Fear God, and give glory to him who sitteth upon the throne, forever and ever; for the hour of his judgment is come.

		% \begin{framed}
		% \scriptsize
			% \begin{compactdesc}
				% \item[JSP] 
				% % \item[BC] 
				% % \item[]
			% \end{compactdesc}
		% \end{framed}

	% -----
	\subsubsection{Section 88:105} % *DC88-105 
	% -----
		\label{DC88-105}

		And again, another angel shall sound his trump, which is the sixth angel, saying: She is fallen who made all nations drink of the wine of the wrath of her fornication; she is fallen, is fallen!

		% \begin{framed}
		% \scriptsize
			% \begin{compactdesc}
				% \item[JSP] 
				% % \item[BC] 
				% % \item[]
			% \end{compactdesc}
		% \end{framed}

	% -----
	\subsubsection{Section 88:106} % *DC88-106 
	% -----
		\label{DC88-106}

		And again, another angel shall sound his trump, which is the seventh angel, saying: It is finished; it is finished! The Lamb of God hath overcome and trodden the wine-press alone,%
			\footnote{%
				% \label{}%
				See \hyperref[ISAIAH-63-3]{Isaiah 63:3} and \hyperref[DC76-107]{DC 76:107}.
				}
		even the wine-press of the fierceness of the wrath of Almighty God.

		% \begin{framed}
		% \scriptsize
			% \begin{compactdesc}
				% \item[JSP] 
				% % \item[BC] 
				% % \item[]
			% \end{compactdesc}
		% \end{framed}

	% -----
	\subsubsection{Section 88:107} % *DC88-107 
	% -----
		\label{DC88-107}

		And then shall the angels be crowned with the glory of his might, and the saints shall be filled with his glory, and receive their inheritance and be made equal with him.

		% \begin{framed}
		% \scriptsize
			% \begin{compactdesc}
				% \item[JSP] 
				% % \item[BC] 
				% % \item[]
			% \end{compactdesc}
		% \end{framed}

	% -----
	\subsubsection{Section 88:108} % *DC88-108 
	% -----
		\label{DC88-108}

		And then shall the first angel again sound his trump in the ears of all living, and reveal the secret acts of men, and the mighty works of God in the first thousand years.

		% \begin{framed}
		% \scriptsize
			% \begin{compactdesc}
				% \item[JSP] 
				% % \item[BC] 
				% % \item[]
			% \end{compactdesc}
		% \end{framed}

	% -----
	\subsubsection{Section 88:109} % *DC88-109 
	% -----
		\label{DC88-109}

		And then shall the second angel sound his trump, and reveal the secret acts of men, and the thoughts and intents of their hearts, and the mighty works of God in the second thousand years---

		% \begin{framed}
		% \scriptsize
			% \begin{compactdesc}
				% \item[JSP] 
				% % \item[BC] 
				% % \item[]
			% \end{compactdesc}
		% \end{framed}

	% -----
	\subsubsection{Section 88:110} % *DC88-110 
	% -----
		\label{DC88-110}

		And so on, until the seventh angel shall sound his trump; and he shall stand forth upon the land and upon the sea, and swear in the name of him who sitteth upon the throne, that there shall be \hyperref[TIME]{time} no longer;%
			\footnote{%
				% \label{}%
				See \hyperref[DC84-100]{DC 84:100}.
				}
		and Satan shall be bound, that old serpent, who is called the devil, and shall not be loosed for the space of a thousand years.

		% \begin{framed}
		% \scriptsize
			% \begin{compactdesc}
				% \item[JSP] 
				% % \item[BC] 
				% % \item[]
			% \end{compactdesc}
		% \end{framed}

	% -----
	\subsubsection{Section 88:111} % *DC88-111 
	% -----
		\label{DC88-111}

		And then he shall be loosed for a little season, that he may gather together his armies.

		% \begin{framed}
		% \scriptsize
			% \begin{compactdesc}
				% \item[JSP] 
				% % \item[BC] 
				% % \item[]
			% \end{compactdesc}
		% \end{framed}

	% -----
	\subsubsection{Section 88:112} % *DC88-112 
	% -----
		\label{DC88-112}

		And Michael, the seventh angel, even the archangel, shall gather together his armies, even the hosts of heaven.

		% \begin{framed}
		% \scriptsize
			% \begin{compactdesc}
				% \item[JSP] 
				% % \item[BC] 
				% % \item[]
			% \end{compactdesc}
		% \end{framed}

	% -----
	\subsubsection{Section 88:113} % *DC88-113 
	% -----
		\label{DC88-113}

		And the devil shall gather together his armies; even the hosts of hell, and shall come up to battle against Michael and his armies.

		% \begin{framed}
		% \scriptsize
			% \begin{compactdesc}
				% \item[JSP] 
				% % \item[BC] 
				% % \item[]
			% \end{compactdesc}
		% \end{framed}

	% -----
	\subsubsection{Section 88:114} % *DC88-114 
	% -----
		\label{DC88-114}

		And then cometh the battle of the great God; and the devil and his armies shall be cast away into their own place, that they shall not have power over the saints any more at all.

		% \begin{framed}
		% \scriptsize
			% \begin{compactdesc}
				% \item[JSP] 
				% % \item[BC] 
				% % \item[]
			% \end{compactdesc}
		% \end{framed}

	% -----
	\subsubsection{Section 88:115} % *DC88-115 
	% -----
		\label{DC88-115}

		For Michael shall fight their battles, and shall overcome him who seeketh the throne of him who sitteth upon the throne, even the Lamb.

		% \begin{framed}
		% \scriptsize
			% \begin{compactdesc}
				% \item[JSP] 
				% % \item[BC] 
				% % \item[]
			% \end{compactdesc}
		% \end{framed}

	% -----
	\subsubsection{Section 88:116} % *DC88-116 
	% -----
		\label{DC88-116}

		This is the glory of God, and the sanctified; and they shall not any more see death.

		% \begin{framed}
		% \scriptsize
			% \begin{compactdesc}
				% \item[JSP] 
				% % \item[BC] 
				% % \item[]
			% \end{compactdesc}
		% \end{framed}

	% -----
	\subsubsection{Section 88:117} % *DC88-117 
	% -----
		\label{DC88-117}

		Therefore, verily I say unto you, my friends, call your solemn assembly, as I have commanded you.

		% \begin{framed}
		% \scriptsize
			% \begin{compactdesc}
				% \item[JSP] 
				% % \item[BC] 
				% % \item[]
			% \end{compactdesc}
		% \end{framed}

	% -----
	\subsubsection{Section 88:118} % *DC88-118 
	% -----
		\label{DC88-118}

		And as all have not faith, seek ye diligently and teach one another words of wisdom; yea, seek ye out of the best books words of wisdom; seek learning, even by study and also by faith.

		% \begin{framed}
		% \scriptsize
			% \begin{compactdesc}
				% \item[JSP] 
				% % \item[BC] 
				% % \item[]
			% \end{compactdesc}
		% \end{framed}

	% -----
	\subsubsection{Section 88:119} % *DC88-119 
	% -----
		\label{DC88-119}

		Organize yourselves; prepare every needful thing; and establish a house, even a house of prayer, a house of fasting, a house of faith, a house of learning, a house of glory, a house of order, a house of God;

		% \begin{framed}
		% \scriptsize
			% \begin{compactdesc}
				% \item[JSP] 
				% % \item[BC] 
				% % \item[]
			% \end{compactdesc}
		% \end{framed}

	% -----
	\subsubsection{Section 88:120} % *DC88-120 
	% -----
		\label{DC88-120}

		That your incomings may be in the name of the Lord; that your outgoings may be in the name of the Lord; that all your salutations may be in the name of the Lord, with uplifted hands unto the Most High.

		% \begin{framed}
		% \scriptsize
			% \begin{compactdesc}
				% \item[JSP] 
				% % \item[BC] 
				% % \item[]
			% \end{compactdesc}
		% \end{framed}

	% -----
	\subsubsection{Section 88:121} % *DC88-121 
	% -----
		\label{DC88-121}

		Therefore, cease from all your light speeches, from all laughter, from all your lustful desires, from all your pride and light-mindedness, and from all your wicked doings.

		% \begin{framed}
		% \scriptsize
			% \begin{compactdesc}
				% \item[JSP] 
				% % \item[BC] 
				% % \item[]
			% \end{compactdesc}
		% \end{framed}

	% -----
	\subsubsection{Section 88:122} % *DC88-122 
	% -----
		\label{DC88-122}

		Appoint among yourselves a teacher, and let not all be spokesmen at once; but let one speak at a time and let all listen unto his sayings, that when all have spoken that all may be edified of all, and that every man may have an equal privilege.

		% \begin{framed}
		% \scriptsize
			% \begin{compactdesc}
				% \item[JSP] 
				% % \item[BC] 
				% % \item[]
			% \end{compactdesc}
		% \end{framed}

	% -----
	\subsubsection{Section 88:123} % *DC88-123 
	% -----
		\label{DC88-123}

		See that ye love one another; cease to be covetous; learn to impart one to another as the gospel requires.

		% \begin{framed}
		% \scriptsize
			% \begin{compactdesc}
				% \item[JSP] 
				% % \item[BC] 
				% % \item[]
			% \end{compactdesc}
		% \end{framed}

	% -----
	\subsubsection{Section 88:124} % *DC88-124 
	% -----
		\label{DC88-124}

		Cease to be idle; cease to be unclean; cease to find fault one with another; cease to sleep longer than is needful; retire to thy bed early, that ye may not be weary; arise early, that your bodies and your minds may be invigorated.

		% \begin{framed}
		% \scriptsize
			% \begin{compactdesc}
				% \item[JSP] 
				% % \item[BC] 
				% % \item[]
			% \end{compactdesc}
		% \end{framed}

	% -----
	\subsubsection{Section 88:125} % *DC88-125 
	% -----
		\label{DC88-125}

		And above all things, clothe yourselves with the bond of charity,%
			\footnote{%
				% \label{}%
				See \hyperref[COLOSSIANS-3-14]{Colossians 3:14}.
				}
		as with a mantle, which is the bond of perfectness and peace.

		% \begin{framed}
		% \scriptsize
			% \begin{compactdesc}
				% \item[JSP] 
				% % \item[BC] 
				% % \item[]
			% \end{compactdesc}
		% \end{framed}

	% -----
	\subsubsection{Section 88:126} % *DC88-126 
	% -----
		\label{DC88-126}

		Pray always, that ye may not faint, until I come. Behold, and lo, I will come quickly, and receive you unto myself. Amen.

		% \begin{framed}
		% \scriptsize
			% \begin{compactdesc}
				% \item[JSP] 
				% % \item[BC] 
				% % \item[]
			% \end{compactdesc}
		% \end{framed}

	% -----
	\subsubsection{Section 88:127} % *DC88-127 
	% -----
		\label{DC88-127}

		And again, the order of the house prepared for the presidency of the school of the prophets, established for their instruction in all things that are expedient for them, even for all the officers of the church, or in other words, those who are called to the ministry in the church, beginning at the high priests, even down to the deacons---

		% \begin{framed}
		% \scriptsize
			% \begin{compactdesc}
				% \item[JSP] 
				% % \item[BC] 
				% % \item[]
			% \end{compactdesc}
		% \end{framed}

	% -----
	\subsubsection{Section 88:128} % *DC88-128 
	% -----
		\label{DC88-128}

		And this shall be the order of the house of the presidency of the school: He that is appointed to be president, or teacher, shall be found standing in his place, in the house which shall be prepared for him.

		% \begin{framed}
		% \scriptsize
			% \begin{compactdesc}
				% \item[JSP] 
				% % \item[BC] 
				% % \item[]
			% \end{compactdesc}
		% \end{framed}

	% -----
	\subsubsection{Section 88:129} % *DC88-129 
	% -----
		\label{DC88-129}

		Therefore, he shall be first in the house of God, in a place that the congregation in the house may hear his words carefully and distinctly, not with loud speech.

		% \begin{framed}
		% \scriptsize
			% \begin{compactdesc}
				% \item[JSP] 
				% % \item[BC] 
				% % \item[]
			% \end{compactdesc}
		% \end{framed}

	% -----
	\subsubsection{Section 88:130} % *DC88-130 
	% -----
		\label{DC88-130}

		And when he cometh into the house of God, for he should be first in the house---behold, this is beautiful, that he may be an example---

		% \begin{framed}
		% \scriptsize
			% \begin{compactdesc}
				% \item[JSP] 
				% % \item[BC] 
				% % \item[]
			% \end{compactdesc}
		% \end{framed}

	% -----
	\subsubsection{Section 88:131} % *DC88-131 
	% -----
		\label{DC88-131}

		Let him offer himself in prayer upon his knees before God, in token or remembrance of the everlasting covenant.

		% \begin{framed}
		% \scriptsize
			% \begin{compactdesc}
				% \item[JSP] 
				% % \item[BC] 
				% % \item[]
			% \end{compactdesc}
		% \end{framed}

	% -----
	\subsubsection{Section 88:132} % *DC88-132 
	% -----
		\label{DC88-132}

		And when any shall come in after him, let the teacher arise, and, with uplifted hands to heaven, yea, even directly, salute his brother or brethren with these words:

		% \begin{framed}
		% \scriptsize
			% \begin{compactdesc}
				% \item[JSP] 
				% % \item[BC] 
				% % \item[]
			% \end{compactdesc}
		% \end{framed}

	% -----
	\subsubsection{Section 88:133} % *DC88-133 
	% -----
		\label{DC88-133}

			\footnote{%
				% \label{}%
				This is too early for Masonry in the LDS church, I think, but this reminds me of the ``five points of fellowship'' in previous versions of the endowment. See \hyperref[TEMPLE-ENDOWMENT-ENDOW-PRE-1990-VEIL]{this part} of the ceremony.
				}%
		Art thou a brother or brethren? I salute you in the name of the Lord Jesus Christ, in token or remembrance of the everlasting covenant, in which covenant I receive you to fellowship, in a determination that is fixed, immovable, and unchangeable, to be your friend and brother through the grace of God in the bonds of love, to walk in all the commandments of God blameless, in thanksgiving, forever and ever. Amen.

		% \begin{framed}
		% \scriptsize
			% \begin{compactdesc}
				% \item[JSP] 
				% % \item[BC] 
				% % \item[]
			% \end{compactdesc}
		% \end{framed}

	% -----
	\subsubsection{Section 88:134} % *DC88-134 
	% -----
		\label{DC88-134}

		And he that is found unworthy of this salutation shall not have place among you; for ye shall not suffer that mine house shall be polluted by him.

		% \begin{framed}
		% \scriptsize
			% \begin{compactdesc}
				% \item[JSP] 
				% % \item[BC] 
				% % \item[]
			% \end{compactdesc}
		% \end{framed}

	% -----
	\subsubsection{Section 88:135} % *DC88-135 
	% -----
		\label{DC88-135}

		And he that cometh in and is faithful before me, and is a brother, or if they be brethren, they shall salute the president or teacher with uplifted hands to heaven, with this same prayer and covenant, or by saying Amen, in token of the same.

		% \begin{framed}
		% \scriptsize
			% \begin{compactdesc}
				% \item[JSP] 
				% % \item[BC] 
				% % \item[]
			% \end{compactdesc}
		% \end{framed}

	% -----
	\subsubsection{Section 88:136} % *DC88-136 
	% -----
		\label{DC88-136}

		Behold, verily, I say unto you, this is an ensample unto you for a salutation to one another in the house of God, in the school of the prophets.

		% \begin{framed}
		% \scriptsize
			% \begin{compactdesc}
				% \item[JSP] 
				% % \item[BC] 
				% % \item[]
			% \end{compactdesc}
		% \end{framed}

	% -----
	\subsubsection{Section 88:137} % *DC88-137 
	% -----
		\label{DC88-137}

		And ye are called to do this by prayer and thanksgiving, as the Spirit shall give utterance in all your doings in the house of the Lord, in the school of the prophets, that it may become a sanctuary, a tabernacle of the Holy Spirit to your edification.

		% \begin{framed}
		% \scriptsize
			% \begin{compactdesc}
				% \item[JSP] 
				% % \item[BC] 
				% % \item[]
			% \end{compactdesc}
		% \end{framed}

	% -----
	\subsubsection{Section 88:138} % *DC88-138 
	% -----
		\label{DC88-138}

		And ye shall not receive any among you into this school save he is clean from the blood of this generation;

		% \begin{framed}
		% \scriptsize
			% \begin{compactdesc}
				% \item[JSP] 
				% % \item[BC] 
				% % \item[]
			% \end{compactdesc}
		% \end{framed}

	% -----
	\subsubsection{Section 88:139} % *DC88-139 
	% -----
		\label{DC88-139}

		And he shall be received by the ordinance of the washing of feet, for unto this end was the ordinance of the washing of feet instituted.%
			\footnote{%
				% \label{}%
				This verse is saying, the ordinance of feet-washing was instituted in order to cleanse people from the ``blood of this generation.'' Think about the \hyperref[TEMPLE-ENDOWMENT-INIT-PRE-2005]{initiatory} here.
				}

		% \begin{framed}
		% \scriptsize
			% \begin{compactdesc}
				% \item[JSP] 
				% % \item[BC] 
				% % \item[]
			% \end{compactdesc}
		% \end{framed}

	% -----
	\subsubsection{Section 88:140} % *DC88-140 
	% -----
		\label{DC88-140}

		And again, the ordinance of washing feet is to be administered by the president, or presiding elder of the church.

		% \begin{framed}
		% \scriptsize
			% \begin{compactdesc}
				% \item[JSP] 
				% % \item[BC] 
				% % \item[]
			% \end{compactdesc}
		% \end{framed}

	% -----
	\subsubsection{Section 88:141} % *DC88-141 
	% -----
		\label{DC88-141}

		It is to be commenced with prayer; and after partaking of bread and wine, he is to gird himself according to the pattern given in the thirteenth chapter of John's testimony concerning me.%
			\footnote{%
				% \label{}%
				See \hyperref[JOHN-13]{John 13}.
				}
		Amen.

		% \begin{framed}
		% \scriptsize
			% \begin{compactdesc}
				% \item[JSP] 
				% % \item[BC] 
				% % \item[]
			% \end{compactdesc}
		% \end{framed}